\scanpage{10}
\mspage{(8)}
Mais quelle est \inserted{(pour $q$ fixé)} la structure de $k\{T\}$-module de $H^{\bullet,q}_{\mathrm{DRpd}}(\mathcal X_*,k)$ ? \deleted{On voit que} pour le \uline{degré total} (égal au degré \unsure{extérieur} : \uline{plus} $q$) on trouve $H^0(\mathcal X_*,k)\otimes_k k\{T\}$ tronqué en degrés $\geq q$, donc
\[
H^{\bullet,q}_{\mathrm{DRpd}}(\mathcal X_*,k)\simeq\tau_q\left(H^0(\mathcal X_*,k)\otimes_k k\{T\}\right)[q]
\]
\margin{(translation des degrés par $-q$)}

Si on réindexe les bidegrés par le couple (degré total, degré extérieur)
\[
H'^{n,q}_{\mathrm{DRpd}}(\mathcal X_*,k)=H^{\overbrace{n-q}^{p},q}_{\mathrm{DRpd}}(\mathcal X_*,k)
\]
(donc la condition de degré $p,q\geq0$ devient $n\geq q\geq0$, l'op. diff est de bidegré $(0,1)$ donc c'est un hom (\uline{homogène}) de $S$-modules gradués), on trouve
\[
H'^{\bullet,q}_{\mathrm{DRpd}}(\mathcal X_*,k)\simeq\tau_q\left(H^q(\mathcal X_*,k)\otimes_k k\{T\}\right)
\]
Ces isom. sont compatibles avec les structures multiplicatives (et bien entendu fonctoriels en $\mathcal X_*,k$...). Donc \inserted{\uline{a priori}} on ne récupère pas \inserted{(via $H'^{\bullet\bullet}_{\mathrm{DRpd}}(\mathcal X_*,k)$)} les $k$-modules $H^q(\mathcal X_*,k)$ et leurs \inserted{cup-}accouplements, mais les $k$-modules \inserted{gradués} $\tau_q(H^q(\mathcal X_*,k)\otimes_k k\{T\})$ et leurs accouplements déduits du cup des $H^\bullet(\mathcal X_*,k)$.
\mspage{9}
